\answer{ex:09:1}{(أ)~$P_\infty=0.2$، والذاكرة $20\,\text{s}$. (ب)~$P_\infty=1$، والذاكرة $1\,\text{s}$. (ج)~الذاكرة تتناسب مع سعة الضجيج: ست مرات.}
\answer{ex:09:2}{$P(t)=P_\infty\coth(t\sqrt{q/R})$. (أ)~$\frac12\sqrt{R/q}=5\,\text{s}$، أي نصف الذاكرة. (ب)~$t=\frac12\ln21\,\sqrt{R/q}\approx15\,\text{s}$: هذه هي المدة التي يجب أن يبقى فيها \nt{ACET} مرتفعًا.}
\answer{ex:09:3}{التباين $qT/2+R/(2T)$؛ $T^*=\sqrt{R/q}$، والتباين $\sqrt{qR}=P_\infty$، كما في المرشح~\eqref{eq:kalman}؛ ويزداد $(\kappa+1/\kappa)/2$ مرة: $1.25$ و$5.05$.}
\answer{ex:09:4}{$a>a_w=1-\theta/(mw)$: $0.15$ عند $w=0.041$، و$0.22$ عند $w=0.045$.}
\answer{ex:09:5}{تظهر الأخطاء عند $M\approx40\text{--}60$، وعند $M=80$ يوجد $1.9$ عصبون زائد، وعند $M=100$ تكون نسبة النمط $0.86$؛ تداخل الأنماط الأخرى تباينه $\approx(M-1)p(1-p)/N$، فيكون $M_{\max}\approx80$، و$M_{\max}\propto N/p$.}
\answer{ex:09:6}{مثلًا $\eta(a)=\eta\min\bigl(1,\max(0,(a-0.3)/0.6)\bigr)$. مع $\eta a$ دخلت العصبونات الغريبة الثلاثة كلها النمط؛ ومع $\eta(a)$ لم يدخل أي منها، والكتابة عند $a\ge0.9$ كما هي (النسبة $1.00$).}
\answer{ex:09:7}{(أ)~$40$ إعادة تشغيل؛ ومع $\eta(a)$ من التمرين~\ref{ex:09:6} أبدًا. (ب)~$200$ إعادة تشغيل، أي $10$ ليالٍ.}
